% !TEX TS-program = xelatex
% !TEX encoding = UTF-8 Unicode
% !Mode:: "TeX:UTF-8"

\documentclass{resume}
\usepackage{zh_CN-Adobefonts_external} % Simplified Chinese Support using external fonts (./fonts/zh_CN-Adobe/)
% \usepackage{zh_CN-Adobefonts_internal} % Simplified Chinese Support using system fonts
\usepackage{linespacing_fix} % disable extra space before next section
\usepackage{cite}

\begin{document}
\pagenumbering{gobble} % suppress displaying page number

\name{王明涵}

% {E-mail}{mobilephone}{homepage}
% be careful of _ in emaill address
\contactInfo{(+86) 136-8325-1930}{wangmh1994@gmail.com}{算法工程师}{GitHub @yuriak}
% {E-mail}{mobilephone}
% keep the last empty braces!
%\contactInfo{xxx@yuanbin.me}{(+86) 131-221-87xxx}{}
 
% \section{个人总结}
% 本人乐于钻研，工作勤奋，热爱尝试挑战性的工作，痴迷于机器学习算法的研究
% 本人在校成绩优秀、乐观向上，工作负责、自我驱动力强、热爱尝试新事物，认同开放、连接、共享的Web在未来的不可替代性。在校期间长期从事可视分析(Web的2D/3D时空可视化)相关研究，对Web技术发展趋势及前端工程化解决方案有浓厚兴趣。\textbf{现任职于华为2012实验室。}

% \section{\faGraduationCap\ 教育背景}
\section{工作经历}
\datedsubsection{\textbf{华为2012实验室\ 翻译中心\ 文本机器翻译实验室}, 算法工程师}{2020.1-Present}
进行研究创新，发表论文，参加国内外算法比赛，以提升团队影响力和技术水准。参与华为自研机器翻译引擎（华为翻译）的开发，用于嵌入鸿蒙OS和HMS Core。
\begin{itemize}
  \item 完成OCR翻译中OCR算法的研发，达到业界先进水平 (ACC>93.7\%)，并集成于华为翻译产品
  \item 完成APE模型从设计、训练到落地，支撑机器翻译交付能力建设
  \item 参加两届WMT，IWSLT比赛，在多个赛道获得第一、二名的成绩
  \item 提升团队科研能力，发表多篇paper
  \item 团队中主持完成了非自回归机器翻译和多模态翻译相关的研究
\end{itemize}

\datedsubsection{\textbf{华为2012实验室\ RAMS实验室}, 算法工程师}{2019.1-2020.1}
构建公司内部使用的轻量级智能客服系统，打通跨部门的资料检索途径，提高部门间信息访问效率。
\begin{itemize}
  \item 设计了低资源场景下的意图识别和槽填充算法，使得系统可以在有限训练数据下做到较高的答准率
  \item 优化了原有系统中NLP相关模块的性能和效率
  \item 重构两套客服系统(IR-based和KG-based)，统一了其中意图识别和槽填充相关的模块
\end{itemize}


\section{教育背景}

\datedsubsection{\textbf{墨尔本大学},计算机科学,\textit{理学硕士学位}}{2017.1 - 2019.1}
\begin{itemize}
    \item GPA:3.75/4.0, \textbf{排名前10\% with distinction}
    \item 相关课程: 机器学习, NLP, AI and Planing, 分布式系统, 统计数据分析, 云计算, 函数式编程等
    \item 毕业论文: Learning Algorithms in Quantitative Trading (supervised by \textbf{Trevor Cohn})
    
    % \item 研究方向：将强化学习和NLP应用于量化交易，使用一个特别设计的RNN模型作为agent通过策略梯度算法学习对投资组合进行调仓，在模型中进一步引入新闻标题后可在VIX期货交易回测中获得70+\%的年化收益
\end{itemize}

\datedsubsection{\textbf{东北财经大学},信息管理与信息系统,\textit{管理学学士学位}}{2012.9 - 2016.7}
\begin{itemize}
    \item GPA:3.7/4.0, \textbf{排名前15\%}
    \item 相关课程: C/C++程序设计, Java程序设计, UML建模, 计算机组成原理, 程序设计方法等
    \item 毕业论文: Online Evaluation System的设计与开发 (supervised by \textbf{李红})
    \item 担任国家教育基金项目研究助理，主导开发了两套教学辅助系统，用于帮助教师提高授课效率并收集课堂数据
    \item 担任大数据研究组研究助理，参与了使用自然语言处理对社交网络数据进行文本分析的研究，使用SVM, Naive Bayes等模型进行舆情监控实验
\end{itemize}

\section{项目经历}
\datedsubsection{\textbf{RLQuant}}{2018.4 - 2019.1}
\begin{itemize}
    \item 将强化学习应用于量化交易的研究 \footnote{https://github.com/yuriak/RLQuant} (研究生毕业论文项目)
    \item 模型通过在回测中探索从而学会做出正确的交易决策，最大化长期收益
    \item 在优化agent的网络结构并加入一个有监督任务稳定训练后，性能可超越原始策略梯度算法10\%
    \item 在模型中进一步引入新闻标题后可在VIX期货交易回测中获得70+\%的年化收益
\end{itemize}

\datedsubsection{\textbf{DLQuant}}{2018.4 - 2019.1}
\begin{itemize}
    \item 将深度学习和自然语言处理应用于量化交易的研究 \footnote{https://github.com/yuriak/DLQuant} (研究生毕业论文项目)
    \item 使用路透社新闻进行美股及商品期货日级别价格变化趋势预测以及使用新浪财经新闻预测A股趋势
    \item 模型使用attention过滤与价格波动相关性更强的新闻, 并通过LSTM编码历史新闻时序关系做出预测
    \item 最终的预测效果与当时SOTA模型持平，但拥有更好的可解释性
\end{itemize}

\datedsubsection{\textbf{自行车赛事赛场管理系统}}{2016.8 - 2016.11}
\begin{itemize}
    \item 深圳唯动体育全栈工程师实习
    \item 独立开发了整套赛场信息服务系统(CISS)
    \item 包括一套后端赛事信息调度平台, 以及一个应用于各赛场手持RFID扫描设备的安卓APP
    \item 该系统已多次稳定运行于多于1500名队员参加的自行车赛事中
\end{itemize}


\section{研究能力}
% increase linespacing [parsep=0.5ex]
\begin{itemize}[parsep=0.2ex]
  \item \textbf{语言与框架}: Python, Shell, Java, C/C++, C\#, LaTeX, PyTorch, Tensorflow 1/2, Mindspore
  \item \textbf{研究兴趣}: 非自回归翻译/生成，强化学习，多模态相关任务，对话系统，量化交易
\end{itemize}

\section{所获荣誉}
\begin{itemize}[parsep=0.2ex]
  \item 华为Mindspore算法挑战赛第一名
  \item 华为公司2020年质量之星奖
  \item 华为北京研究所所长奖 --- 杰出贡献专家奖
  \item IWSLT 2021 Offline Speech Translation 第一名
  \item IWSLT 2020 Video Speech Translation 第一名
  \item WMT 2021 QE, Domain, News 等多个赛道第一名
  \item WMT 2020 QE, APE, Domain, News 等多个赛道第一名
  \item SemEval-2019 地名解析比赛第二名
  \item 2014年花旗杯金融信息科技竞赛优胜奖
\end{itemize}

\section{研究成果}
\begin{itemize}[parsep=0.2ex]
  \item How does Length Prediction Influence the Performance of Non-Autoregressive Translation? (BlackboxNLP 2021)
  \item HI-CMLM: Improve CMLM with Hybrid Decoder Input (INLG 2021)
  \item Make the Blind Translator See The World: A Novel Transfer Learning Solution for Multimodal Machine Translation (AMTA 2021)
  \item Efficient Transfer Learning for Quality Estimation with Bottleneck Adapter Layer (EAMT 2020)
  \item Unified Humor Detection Based on Sentence-pair Augmentation and Transfer Learning (EAMT 2020)
  \item SMAC: An Interpretable Reasoning Network for Visual Question Answering (RCQA 2020)
  \item The HW-TSC's Offline Speech Translation Systems for IWSLT 2021 Evaluation (IWSLT 2021)
  \item The HW-TSC Video Speech Translation System at IWSLT 2020 (IWSLT 2020)
  \item HW-TSC's Participation at WMT 2020 Quality Estimation Shared Task (WMT 2020)
  \item HW-TSC's Participation at WMT 2020 Automatic Post Editing Shared Task (WMT 2020)
  \item Huawei's Submissions to the WMT20 Biomedical Translation Task (WMT 2020)
  \item HW-TSC's Participation in the WMT 2020 News Translation Shared Task (WMT 2020)
  \item UniMelb at SemEval-2019 Task 12: Multi-model Combination for Toponym Resolution (SemEval 2019)
\end{itemize}

% \ GPA:3.7/4, \textbf{排名前15\%}
% 毕业论文: Online Evaluation System的设计与开发

% \section{研究能力}
% % increase linespacing [parsep=0.5ex]
% \begin{itemize}[parsep=0.2ex]
%   \item \textbf{语言与框架}: Python, Shell, Java, C/C++, C\#, LaTeX, PyTorch, Tensorflow 1/2, Mindspore
%   \item \textbf{研究兴趣}: 非自回归翻译/生成，强化学习，多模态相关任务，对话系统，量化交易
% \end{itemize}

% \end{itemize}

% \section{实习经历}
% \datedsubsection{\textbf{阿里巴巴集团 | Alibaba}, 前端开发工程师}{2017.6-2017.9}
% \begin{itemize}
% %   \item 飞猪北京前端团队全面负责各交通线的票务(机票/火车票/汽车票) web 应用与事业群基础架构研发
%   \item 独立负责车站地图开发(React),通过HTML5 本地存储及JSBridge实现在阿里全系应用中发布上线
%   \item 独立负责BU SPM chrome插件开发,支付成功/订单详情等页面的开发与交叉营销的接入工作
% \end{itemize}

% \datedsubsection{\textbf{北京腾云天下科技有限公司 | TalkingData},数据挖掘与可视化工程师}{2015.11-2017.5}
% \begin{itemize}
%   \item \textbf{利用海量用户定位数据，对城市空间及人群移动特征进行研究。}第一个课题是基于香农熵和人群出行模式，构建城市网格与用户矩阵分析城市多样性/流动性分布；可视分析平台前端与可视化基于D3/Vue/Express开发，数据分析与存储采用Python/MySQL/MongoDB技术，为了均衡大数据情况下的页面可视化渲染消耗用canvas替代svg。第二个课题是对海量商场定位数据做人群分类与可视化查询，依据该系统撰写的论文被CIKM 2016(DAVA Workshop)录用，并收录于中科院软件所年会成果集
%   \item 负责数据科学部HQ LAB的可视化原型开发，主导 TalkingMind 平台系统设计与前端开发
% \end{itemize}

% \datedsubsection{\textbf{北京格灵深瞳信息技术有限公司 | DeepGlint},Web开发工程师}{2015.7-2015.9}
% \begin{itemize}
%   \item \textbf{独立负责MUSE部门的可视化组件研发。}与平台研发、设计协作完成 DeepGlint Developer 平台可视化图表组件的集成开发，符合完全定制化渲染、响应式布局与实时更新等特点
%   \item 利用 D3+Vue+WebGL(Three.js) 尝试实现三维空间的人群移动可视化
% \end{itemize}

% % \begin{onehalfspacing}
% % \end{onehalfspacing}

% % \datedsubsection{\textbf{DID-ACTE} 荷兰莱顿}{2015年3月 - 2015年6月}
% % \role{本科毕业设计}{LIACS 交换生}
% % 利用结巴分词对中国古文进行分词与词性标注，用已有领域知识训练形成 classifier 并对结果进行调优
% % \begin{onehalfspacing}
% % \begin{itemize}
% %   \item 利用结巴分词对中国古文进行分词与词性标注
% %   \item 利用已有领域知识训练形成 classifier, 并用分词结果进行测试反馈
% %   \item 尝试不同规则，对 classifier 进行调优
% % \end{itemize}
% % \end{onehalfspacing}

% \section{竞赛获奖/项目作品}
% % increase linespacing [parsep=0.5ex]
% \begin{itemize}[parsep=0.2ex]
% %   \item LeetCodeOJ Solutions, \textit{https://github.com/hijiangtao/LeetCodeOJ}
%   \item 第三届中国软件杯大学生软件设计大赛\textbf{全国一等奖}( \textit{http://www.cnsoftbei.com/} ),2014 年8月
%   \item 中国机器人大赛创意设计大赛\textbf{全国特等奖}( \textit{http://www.rcccaa.org/} ),2013年8月
% %   \item 中国机器人大赛暨Robocup公开赛(武术擂台赛)全国一等奖,2013年10月
%   \item 第11届北京理工大学“世纪杯”竞赛学生课外科技作品竞赛\textbf{特等奖},2013年8月
%   \item VIS Components for security system, \textit{https://hijiangtao.github.io/ss-vis-component/}
%   \item 个人博客：\textit{https://hijiangtao.github.io/}，更多作品见 \textit{https://github.com/hijiangtao}
% %   \item 电视节目"爸爸去哪儿"可视化分析展示, \textit{https://hijiangtao.github.io/variety-show-hot-spot-vis/}
% \end{itemize}

% % \section{\faHeartO\ 项目/作品摘要}
% % \section{项目/作品摘要}
% % \datedline{\textit{An Integrated Version of Security Monitor Vis System}, https://hijiangtao.github.io/ss-vis-component/ }{}
% % \datedline{\textit{Dark-Tech}, https://github.com/hijiangtao/dark-tech/ }{}
% % \datedline{\textit{融合社交网络数据挖掘的电视节目可视化分析系统}, https://hijiangtao.github.io/variety-show-hot-spot-vis/}{}
% % \datedline{\textit{LeetCodeOJ Solutions}, https://github.com/hijiangtao/LeetCodeOJ}{}
% % \datedline{\textit{Info-Vis}, https://github.com/ISCAS-VIS/infovis-ucas}{}


% % \section{\faInfo\ 社会实践/其他}
% \section{社区参与/实践其他}
% % increase linespacing [parsep=0.5ex]
% \begin{itemize}[parsep=0.2ex]
%   \item 乐于参与开源社区讨论,\textbf{参与翻译 Vue.js, webpack, WebAssembly, Babel 文档，印记中文成员}
%   \item 中国科学院大学2016秋季学期可视化与可视分析课程助教，\textit{http://vis.ios.ac.cn/infovis-ucas/}
%   \item 未来论坛学生会成员、北理社联新闻信息中心主任、北理工软件学院学生会宣传部副部长(2012-2016)
%   \item 2013-2015 北京市共青团“温暖衣冬”志愿者,第九届园博会志愿者,2014 FLL机器人世锦赛志愿者
% \end{itemize}

%% Reference
%\newpage
%\bibliographystyle{IEEETran}
%\bibliography{mycite}
\end{document}
